\documentclass{article}

\usepackage[utf8]{inputenc}
\usepackage[portuguese]{babel}
\usepackage[T1]{fontenc}
\usepackage{lmodern}
\usepackage{fullpage}
\usepackage[usenames]{color}
\usepackage[color]{coqdoc}
\usepackage{url}
\usepackage{makeidx,hyperref}
\usepackage[all]{xy}

\usepackage{mathpartir}
\usepackage{tcolorbox}
\usepackage{amsmath,amssymb}

\newcommand{\tto}{\twoheadrightarrow}
\newcommand{\ott}{\twoheadleftarrow}

\title{Formalização da correção do algoritmo de ordenação por inserção binária}
\author{
  Daniel da Cunha Pereira Luz \textsc -- 211055540 \\
  Thiago Veras Rodrigues Queiroz \textsc -- 211055370 \\
  Vinícius Bowen \textsc -- 180079239
}
\date{\today}

\begin{document}
\maketitle

\section{Introdução}

Este relatório documenta a formalização, no assistente de provas Rocq, da
correção do algoritmo de ordenação por inserção binária
(\emph{binary insertion sort}). O algoritmo, implementado pela função
\texttt{binsertion\_sort}, insere cada elemento de uma lista em sua posição
correta dentro de uma sublista já ordenada, usando uma busca binária
(\texttt{bsearch}) para localizar essa posição em tempo logarítmico, em vez
da busca linear usada na inserção tradicional.

Formalmente, a correção do algoritmo é estabelecida por dois resultados
combinados: que a lista produzida por \texttt{binsertion\_sort} está
ordenada (\texttt{Sorted le}) e que ela é uma permutação da lista de
entrada (\texttt{Permutation}). O código-fonte completo com todas as
definições e provas está disponível no repositório do projeto; neste
documento, cada integrante descreve em linguagem natural a estratégia de
prova da parte sob sua responsabilidade.

\section{Fundamentos e Propriedades Estruturais}
\label{sec:daniel}

Esta seção detalha o desenvolvimento dos blocos de construção fundamentais do projeto, sob a responsabilidade do integrante Daniel, localizados no arquivo \texttt{src/binsertion\_sort.v}. Esses blocos garantem a validade matemática dos limites das buscas e descrevem as propriedades estruturais de manipulação de listas ao realizar uma inserção direta.

\subsection{Testes de Execução Computacional}

Antes de iniciar as demonstrações de teoremas gerais, foram definidos cinco lemas de testes computacionais (\coqdoclemma{test0} a \coqdoclemma{test4}). Esses testes realizam a avaliação simbólica e a redução de expressões concretas no interpretador do Rocq. A estratégia de prova desses testes baseia-se na reescrita da equação recursiva da busca binária (\coqdockw{rewrite} \coqdoclemma{bsearch\_equation}) seguida pela simplificação direta (\coqdoctac{simpl}) e reflexividade de igualdade (\coqdoctac{reflexivity}). No caso de barreira onde elementos precisam ser comparados aritmeticamente, utilizou-se a inversão de hipóteses absurdas (\coqdoctac{inversion}) para descartar caminhos impossíveis. Esses testes cumprem a função de sanidade do código, garantindo que a especificação de \coqdocdefinition{bsearch} computa os índices esperados na prática.

\subsection{Abordagem de Prova: Limites da Busca (\texttt{bsearch\_valid\_pos})}

A função \coqdocdefinition{bsearch} calcula o índice ideal de inserção dividindo recursivamente a lista original ao meio usando as funções auxiliares \coqdocdefinition{firstn} (prefixo) e \coqdocdefinition{skipn} (sufixo). O lema \coqdoclemma{bsearch\_valid\_pos} estabelece que o índice retornado por essa busca sempre respeita os limites físicos da lista, isto é, está no intervalo fechado entre $0$ e o tamanho total da lista.

\begin{mathpar}
\inferrule* [lab={bsearch\_bounds}]
  { \coqdocvariable{x} \in \mathbb{N} \\ \coqdocvariable{l} \in \coqdocinductive{list} \ \mathbb{N} }
  { 0 \le \coqdocdefinition{bsearch} \ \coqdocvariable{x} \ \coqdocvariable{l} \le \text{length}(\coqdocvariable{l}) }
\end{mathpar}

Como a redução do tamanho da lista em \coqdocdefinition{bsearch} não é linear (não decrementa elemento por elemento), a indução estrutural padrão falha. A abordagem de prova adotada baseia-se em \textbf{indução funcional} (\coqdoctac{functional induction}), estruturando a demonstração passo a passo conforme as ramificações de execução do algoritmo:

\begin{enumerate}
  \item \textbf{Caso Base 1 (Lista Vazia):} Se a lista for vazia, a função retorna $0$. Como o comprimento de uma lista vazia é $0$, o limite $0 \le 0 \le 0$ é imediato por aritmética elementar.
  
  \item \textbf{Caso Base 2 (Lista Unitária):} Se a lista contém apenas um elemento $y$, o algoritmo testa se $x \le y$. Se verdadeiro, retorna $0$; se falso, retorna $1$. Em ambos os caminhos de decisão lógica, o índice resultante está contido nos limites do tamanho da lista unitária ($1$), o que é verificado por análise de casos e simplificação.
  
  \item \textbf{Caso Geral 1 (Subdivisão à Esquerda):} Quando o elemento procurado é menor do que o ponto médio, o algoritmo recorre sobre o prefixo da lista obtido via \coqdocdefinition{firstn}. A hipótese de indução garante que o índice retornado é válido dentro dos limites do prefixo. Para estender essa garantia à lista original, a prova reescreve a relação estrutural usando o lema \coqdoclemma{length\_firstn} (que estabelece o comprimento do prefixo como o valor mínimo entre o ponto médio de corte e o tamanho total da lista). Com isso, a propriedade é reduzida a uma desigualdade estritamente aritmética solucionada pelo resolvedor linear (\coqdoctac{lia}).
  
  \item \textbf{Caso Geral 2 (Ponto Médio):} Se o elemento procurado coincide com o ponto médio da lista, a função encerra a recursão e retorna diretamente o índice médio. Como a lista neste caso geral possui comprimento maior ou igual a $2$, a prova do limite $0 \le \text{ponto\_m\text{é}dio} \le \text{comprimento}$ decorre de propriedades de divisão inteira por $2$, auxiliadas pelo teorema \coqdoclemma{Nat.div\_lt}.
  
  \item \textbf{Caso Geral 3 (Subdivisão à Direita):} Quando o elemento procurado é maior do que o ponto médio, o algoritmo recorre sobre o sufixo obtido via \coqdocdefinition{skipn}, retornando a soma do índice médio com o resultado dessa chamada recursiva. A hipótese de indução garante a validade do índice na sublista. A prova do limite global é concluída utilizando o lema \coqdoclemma{length\_skipn} para relacionar o tamanho do sufixo à diferença aritmética do comprimento total. O resolvedor aritmético unifica os termos e fecha o objetivo demonstrando que a soma das partições preserva o limite superior original.
\end{enumerate}

\subsection{Abordagem de Prova: Propriedades Estruturais de \texttt{insert\_at}}

A função \coqdocdefinition{insert\_at} \coqdocvariable{i} \coqdocvariable{x} \coqdocvariable{l} insere o elemento \coqdocvariable{x} no índice \coqdocvariable{i} da lista \coqdocvariable{l}. As demonstrações das propriedades essenciais dessa função seguiram abordagens indutivas bem delimitadas:

\begin{itemize}
  \item \textbf{Prova de Conservação de Tamanho (\coqdoclemma{insert\_at\_length}):} O lema garante que a inserção incrementa o comprimento da lista em exatamente uma unidade.
  
  \begin{mathpar}
  \inferrule* [lab={insert\_length}]
    { \coqdocvariable{x}, \coqdocvariable{i} \in \mathbb{N} \\ \coqdocvariable{l} \in \coqdocinductive{list} \ \mathbb{N} }
    { \text{length}(\coqdocdefinition{insert\_at} \ \coqdocvariable{i} \ \coqdocvariable{x} \ \coqdocvariable{l}) = S(\text{length}(\coqdocvariable{l})) }
  \end{mathpar}
  
  A abordagem utilizou indução estrutural sobre a lista \coqdocvariable{l} combinada com a análise de casos sobre o índice de destino \coqdocvariable{i}. Para o caso base (lista vazia), a inserção resulta trivialmente em uma lista unitária de tamanho $1$. Para o caso indutivo, se o índice for sucessor de outro valor, a recursão percorre a lista até encontrar a posição correta. O passo indutivo foi solucionado aplicando o lema de igualdade de construtores \coqdoctac{f\_equal} (reduzindo a igualdade de sucessores à igualdade de seus argumentos), permitindo o casamento perfeito com a hipótese de indução.

  \item \textbf{Prova de Preservação de Conteúdo (\coqdoclemma{insert\_at\_perm}):} Demonstra que a lista resultante após a inserção é uma permutação válida da lista de entrada com o novo elemento anexado à cabeça (\coqdocvariable{x} :: \coqdocvariable{l}).
  
  A prova seguiu por indução estrutural na lista \coqdocvariable{l}. O caso base e o subcaso de inserção na cabeça da lista ($i=0$) são triviais. No caso indutivo geral, onde o elemento é inserido na cauda recursivamente, a hipótese de indução garante que a sublista alterada é uma permutação de sua forma com $x$ anexado. A partir dessa premissa, a prova aplicou a regra de congruência de cabeça de permutações (\coqdoclemma{perm\_skip}) para estender a relação e utilizou as regras de transitividade (\coqdoclemma{perm\_trans}) e simetria de vizinhos adjacentes (\coqdoclemma{perm\_swap}) da biblioteca padrão do Rocq para provar que a alteração de posição intermediária preserva a equivalência do conjunto final.
\end{itemize}

\subsection{Desafios e Adaptações de Automação no Rocq 9}

Durante o desenvolvimento no Rocq 9, o principal desafio residiu na instabilidade da tática de indução funcional ao lidar com nomes dinâmicos de variáveis gerados pelo compilador (como sublistas intermediárias). Scripts de prova que citavam variáveis diretamente falhavam em atualizações do ambiente do assistente. 

Para contornar essa fragilidade e tornar as demonstrações robustas, o script de \coqdoclemma{bsearch\_valid\_pos} foi reestruturado para ser executado em lote por meio do combinador \coqdoctac{all:}, aplicando táticas baseadas unicamente em estruturas de tipos (como \coqdoclemma{length\_firstn} e \coqdoclemma{length\_skipn}) e permitindo que o resolvedor aritmético resolvesse todos os ramos recursivos simultaneamente sem apelo a nomes dinâmicos. Adicionalmente, na prova de tamanho, a hipótese de indução apresentava-se sob uma quantificação universal aberta (\coqdockw{forall} \coqdocvariable{i} \coqdocvariable{x}), o que impedia sua aplicação direta pelo resolvedor automático. Isso foi solucionado pela introdução da tática de desestruturação funcional \coqdoctac{f\_equal} para remover construtores externos antes da unificação dos tipos.

O desenvolvimento bem-sucedido dessas propriedades forneceu ao Thiago e ao Vinícius as garantias de estabilidade de índices e invariantes de permutações necessárias para a prova dos teoremas de ordenação global do projeto.

\section{Correção de \texttt{binsert}}
\label{sec:thiago}

Nesta seção é apresentada a prova de que a função \texttt{binsert}, responsável
por inserir um elemento \texttt{x} em uma lista ordenada \texttt{l} na posição
retornada por \texttt{bsearch}, preserva a ordenação da lista. Em outras
palavras, mostra-se que, se \texttt{l} está ordenada, então
\texttt{insert\_at (bsearch x l) x l} também estará.

\subsection{Estratégia geral}

A demonstração foi dividida em duas etapas principais.

\begin{enumerate}
  \item Primeiro, foi provado um resultado sobre \texttt{insert\_at},
        independente da função \texttt{bsearch}. A ideia é mostrar que,
        se a posição de inserção separa corretamente a lista --- isto é,
        todos os elementos anteriores são menores ou iguais a
        \texttt{x}, enquanto todos os elementos a partir dessa posição são
        maiores ou iguais a \texttt{x} --- então a inserção preserva a
        ordenação da lista.

  \item Em seguida, foi demonstrado que a posição retornada por
        \texttt{bsearch} satisfaz exatamente essa propriedade sempre que a
        lista de entrada já estiver ordenada.
\end{enumerate}

A combinação desses dois resultados permite concluir diretamente o teorema
principal, uma vez que \texttt{binsert x l} é definido como
\texttt{insert\_at (bsearch x l) x l}.

\subsection{Lemas auxiliares sobre listas ordenadas}

Antes da prova principal, foi necessário estabelecer alguns resultados
auxiliares sobre listas ordenadas que não estavam disponíveis na biblioteca
padrão na forma necessária para esta formalização.

\begin{itemize}
  \item Em uma lista ordenada \texttt{a :: l}, o elemento \texttt{a} é
        menor ou igual a todos os elementos de \texttt{l}, e não apenas ao
        primeiro.

  \item O prefixo (\texttt{firstn}) de uma lista ordenada também permanece
        ordenado.

  \item O sufixo (\texttt{skipn}) de uma lista ordenada também permanece
        ordenado.

  \item Em uma lista escrita como \texttt{P ++ a :: Q}, todos os elementos
        de \texttt{P} são menores ou iguais a \texttt{a}.
\end{itemize}

Embora esses resultados sejam esperados para listas ordenadas, foi
necessário demonstrá-los explicitamente por indução, já que a definição de
\texttt{Sorted} fornece apenas propriedades entre elementos consecutivos.

\subsection{Ligando \texttt{bsearch} à ordenação}

O principal resultado desta etapa mostra que a posição retornada por
\texttt{bsearch} realmente divide a lista em dois grupos: antes da posição
ficam apenas elementos menores ou iguais a \texttt{x}, enquanto a partir
dela ficam apenas elementos maiores ou iguais.

A demonstração foi feita por indução funcional sobre a definição de
\texttt{bsearch}, acompanhando seus casos: lista vazia, lista unitária e
lista dividida ao meio.

O caso de maior complexidade ocorre quando a lista é particionada em duas
metades. Nessa situação, \texttt{bsearch} realiza uma chamada recursiva
sobre apenas uma das metades, e a posição final é obtida somando o
deslocamento correspondente ao ponto médio. A prova consiste em mostrar que
essa posição continua separando corretamente a lista original. Para isso,
foi utilizada a hipótese de indução sobre a metade escolhida em conjunto
com os lemas sobre listas ordenadas apresentados anteriormente.

\subsection{Uma dificuldade encontrada}

Durante o desenvolvimento dessa prova foi necessário alterar a estratégia
utilizada no caso em que a busca continua na metade direita da lista.

Inicialmente, a decomposição da lista era feita por meio de reescritas com
\texttt{replace ... at N}, escolhendo manualmente qual ocorrência deveria
ser substituída. Essa abordagem acabou se mostrando pouco robusta, pois o
uso de \texttt{set} para nomear a lista e o ponto médio fazia com que a
numeração das ocorrências variasse durante a prova, tornando a aplicação da
tática instável.

A solução adotada foi substituir essa abordagem por dois lemas gerais sobre
\texttt{firstn} e \texttt{skipn}, relacionando diretamente índices da forma
\texttt{n + k} com aplicações sucessivas dessas funções. Dessa forma, a
prova deixou de depender da posição das ocorrências na expressão e passou a
utilizar apenas propriedades estruturais das listas, tornando-se mais
simples e estável.

\subsection{Fechando a prova de \texttt{binsert\_correct}}

Após estabelecer o lema de separação e o resultado sobre
\texttt{insert\_at}, a demonstração do teorema
\texttt{binsert\_correct} torna-se direta. Basta aplicar o lema de
separação à posição retornada por \texttt{bsearch} e utilizar
\texttt{bsearch\_valid\_pos} para garantir que essa posição pertence aos
limites da lista.

\end{document}